\section{模的乘积}

记号约定:$A$是环,$\left(E_{i}\right)_{i\in I},\left(F_{i}\right)_{i\in I},\left(G_{i}\right)_{i\in I}$是$A$-模构成的族,$F$是$A$-模.

\begin{definition}{模的乘积}{}
	$\left(E_{i}\right)_{i\in I}$上的模结构使$\prod\limits_{i\in I}E_{i}$是$A$-模,称$\prod\limits_{i\in I}E_{i}$为$\left(E_{i}\right)_{i\in I}$的乘积模.
\end{definition}

\begin{remark}{}{}
	\begin{enumerate}
		\item $\left(\pr_{i}\right)_{i\in I}$是一族$A$-线性映射.
		\item 当$I=\emptyset$时,$\prod\limits_{i\in I}E_{i}$是单点集,进而是零模.
	\end{enumerate}
\end{remark}

\begin{proposition}{模的乘积的泛性质}{}
	记$E=\prod\limits_{i\in I}E_{i}$.设$\forall i\in I\left(f_{i}\in\Hom_{A}\left(F,E_{i}\right)\right)$,则存在唯一的$f\in\Hom_{A}\left(F,E\right)$使下图交换.
	\begin{center}
		\begin{tikzcd}
			F && E \\
			& {E_{i}}
			\arrow["\exists!f", dashed, from=1-1, to=1-3]
			\arrow["{f_{i}}"', from=1-1, to=2-2]
			\arrow["{\pr_{i}}", two heads, from=1-3, to=2-2]
		\end{tikzcd}
	\end{center}
	具体的,$f:F\to E,x\mapsto\left(f_{i}\left(x\right)\right)_{i\in I}$.
\end{proposition}

\begin{proposition}{模的乘积的``结合律''}{}
	设$I\neq\emptyset$且有划分$\left(I_{\lambda}\right)_{\lambda\in\Lambda}$,则有典范同构$\prod\limits_{i\in I}E_{i}\simeq\prod\limits_{\lambda\in\Lambda}\prod\limits_{i\in I_{\lambda}}E_{i}$.
\end{proposition}

\begin{proposition}{线性映射族的乘积是线性映射}{}
	设$\forall i\in I\left(f_{i}\in\Hom_{A}\left(E_{i},F_{i}\right)\right)$,则$f:\prod\limits_{i\in I}E_{i}\to\prod\limits_{i\in I}F_{i},\left(x_{i}\right)_{i\in I}\mapsto\left(f_{i}\left(x_{i}\right)\right)_{i\in I}$是$A$-线性映射.
\end{proposition}

\begin{definition}{线性映射族的乘积}{}
	设$\forall i\in I\left(f_{i}\in\Hom_{A}\left(E_{i},F_{i}\right)\right)$.称$\prod\limits_{i\in I}f_{i}:\prod\limits_{i\in I}E_{i}\to\prod\limits_{i\in I}M_{i},\left(x_{i}\right)_{i\in I}\mapsto \left(f_{i}\left(x_{i}\right)\right)_{i\in I}$为$\left(f_{i}\right)_{i\in I}$的乘积.
\end{definition}

\begin{proposition}{线性映射族的乘积与正合性}{}
	设$\forall i\in I\left(f_{i}\in\Hom_{A}\left(E_{i},F_{i}\right)\wedge g_{i}\in\Hom_{A}\left(F_{i},G_{i}\right)\right)$.记$E=\prod\limits_{i\in I}E_{i},F=\prod\limits_{i\in I}F_{i},G=\prod\limits_{i\in I}G_{i},f=\prod\limits_{i\in I}f_{i},g=\prod\limits_{i\in I}g_{i}$,则
	\begin{center}
		$E\xlongrightarrow{f}F\xlongrightarrow{g}G$正合$\iff$对任一$i\in I$,$E_{i}\xlongrightarrow{f_{i}}F_{i}\xlongrightarrow{g_{i}}G_{i}$正合.
	\end{center}
\end{proposition}

\begin{corollary}{}{}
	设$\forall i\in I\left(f_{i}\in\Hom_{A}\left(E_{i},F_{i}\right)\right)$.记$E=\prod\limits_{i\in I}E_{i},F=\prod\limits_{i\in I}F_{i}$.
	\begin{enumerate}
		\item $\ker\left(f\right)=\prod\limits_{i\in I}\ker\left(f_{i}\right)$,$\im\left(f\right)=\prod\limits_{i\in I}\im\left(f_{i}\right)$.
		\item 存在典范同构$\coim\left(f\right)\simeq\prod\limits_{i\in I}\coim\left(f_{i}\right)$,$\coker\left(f\right)\simeq\prod\limits_{i\in I}\coker\left(f_{i}\right)$.
		\item $f$是单射(相对地,满射,双射)$\iff$ $\left(f_{i}\right)_{i\in I}$是一族单射(相对地,满射,双射).
	\end{enumerate}
\end{corollary}

\begin{proposition}{}{}
	设对任意$i\in I$,$F_{i}$是$E_{i}$的子模.记$E=\prod\limits_{i\in I}E_{i},F=\prod\limits_{i\in I}F_{i}$,则$F$是$E$的子模,并且有典范同构$\prod\limits_{i\in I}\left(E_{i}/F_{i}\right)\simeq E/F$.
\end{proposition}

\begin{definition}{幂模}{}
	我们称$E^{I}\coloneq\prod\limits_{i\in I}E$为$E$的$I$次幂模.
\end{definition}

\begin{proposition}{对角映射是$A$-线性映射}{}
	对角映射$\Delta:E\to E^{I}$是$A$-线性映射.
\end{proposition}

\begin{proposition}{诱导映射的分解}{}
	设$\forall i\in I\left(f_{i}\in\Hom_{A}\left(F,E\right)\right)$,$f:F\to E^{I},x\mapsto\left(f_{i}\left(x\right)\right)_{i\in I}$,则下图交换.
	\begin{center}
		\begin{tikzcd}
			F && {F^{I}} \\
			& {\prod\limits_{i\in I}E_{i}}
			\arrow["\Delta", from=1-1, to=1-3]
			\arrow["f"', from=1-1, to=2-2]
			\arrow["{\prod\limits_{i\in I}f_{i}}", from=1-3, to=2-2]
		\end{tikzcd}
	\end{center}
\end{proposition}


